\begingroup
\captionsetup{hypcap=false}
\captionof{table}{Exclusion Criteria}
\label{tab:exclusion-criteria}
\endgroup
\begin{tcolorbox}[
  enhanced,
  breakable,
  colback=white,
  colframe=excred,
  title={\textbf{Exclusion Criteria}},
  fonttitle=\small\bfseries,
  coltitle=white,
  attach boxed title to top left={yshift=-2mm, xshift=6mm},
  boxed title style={colback=excred, colframe=excred},
  top=6pt, bottom=6pt, left=8pt, right=8pt,
  before upper={\small}
]
\setlength{\tabcolsep}{0pt}%
\begin{tabular}{%
  @{}%
  >{\bfseries\arraybackslash}p{0.9cm}%
  >{\arraybackslash}p{\dimexpr\tcbtextwidth-0.9cm\relax}%
  @{}%
}
EC1 & Sources prior to 2018 unless meeting IC7; all exceptions recorded in the review log. \\[3pt]
EC2 & Non-peer-reviewed opinion pieces, blog posts, or media articles without evidential grounding, unless from authoritative professional or statutory bodies. \\[3pt]
EC3 & Studies focused solely on adult sexual offending with no transferable relevance to child safeguarding contexts. \\[3pt]
EC4 & Technical \gls{nlp} literature without contextual application to safeguarding, criminal justice, or operational deployment (foundational instructional / code text not included)\\[3pt]
EC5 & Technology adoption studies from sectors lacking structural or cultural comparability to UK policing (e.g.\ private-sector contexts without analogous accountability frameworks). \\[3pt]
EC6 & Duplicate publications; where multiple versions exist, the most recent peer-reviewed version is retained. \\[3pt]
EC7 & Sources for which full text cannot be obtained via institutional access, inter-library loan, or open-access repositories within a reasonable retrieval window. \\
\end{tabular}
\end{tcolorbox}